\chapter{ANALYSE SWOT}
\thispagestyle{empty}

\begin{itemize}
    \item \textbf{Forces (Strengths) :}
    \begin{itemize}
        \item La structure d'hyperliens de Wikipédia présente naturellement les relations entre articles sous forme de graphe ; il n'est pas nécessaire de construire un graphe séparément.
        \item Les modèles de plongements de phrases utilisés (bge-m3 et gte-modernbert-base), étant pré-entraînés sur de grands corpus textuels, sont capables de produire des représentations sémantiques de haute qualité sans nécessiter d'adaptation supplémentaire au domaine.
        \item Le mécanisme d'attention du GAT améliore la qualité des recommandations en apprenant automatiquement l'importance de chaque lien.
        \item L'existence de jeux de données de référence prêts à l'emploi et nettoyés tels que Wiki-CS \cite{mernyei2020wikics} permet de tester rapidement le modèle.
    \end{itemize}

    \item \textbf{Faiblesses (Weaknesses) :}
    \begin{itemize}
        \item La limite de longueur d'entrée des modèles de plongements de phrases empêche de traiter l'intégralité des articles Wikipédia longs ; des stratégies de troncature ou de résumé seront nécessaires.
        \item L'exécution conjointe du GNN et des modèles de plongements requiert une grande quantité de mémoire GPU ; cette situation peut créer un goulet d'étranglement en raison des contraintes matérielles.
        \item Le projet nécessite une expertise dans deux domaines distincts : l'ingénierie des données et le développement de modèles.
    \end{itemize}

    \item \textbf{Opportunités (Opportunities) :}
    \begin{itemize}
        \item La modernisation de la couche de représentation textuelle de WikiRecNet \cite{wikirecnet2020}, qui utilise Doc2Vec, avec des modèles modernes de plongements de phrases offre une opportunité directe d'amélioration.
        \item L'utilisation des LLM comme couche de présentation hiérarchique constitue un nouveau domaine de recherche encore peu exploré dans la littérature.
        \item La méthode développée possède le potentiel d'être étendue à d'autres domaines de Wikipédia au-delà de l'informatique, ainsi qu'à différentes langues.
    \end{itemize}

    \item \textbf{Menaces et risques (Threats) :}
    \begin{itemize}
        \item La nature constamment évolutive de Wikipédia comporte le risque que le jeu de données perde sa pertinence au fil du temps.
        \item Des problèmes de mémoire insuffisante (OOM) peuvent survenir lors de l'entraînement du GNN sur un graphe de grande taille.
        \item L'absence d'un jeu de données de vérité terrain (ground truth) approprié pour l'évaluation quantitative de la qualité des recommandations complique le processus d'évaluation.
        \item Le risque d'hallucination de la couche LLM peut affecter négativement la fiabilité de la présentation hiérarchique.
    \end{itemize}
\end{itemize}